\subsection{西罗子群}

定义 10. 设 G 是一个有限群。G 的西罗 p-子群指 G 的满足以下两个条件的任意子群 P：

\begin{enumerate}
\def\labelenumi{(\alph{enumi})}
\item
  P 是一个 p-群。
\item
  (G:P) 不是 p 的倍数。
\end{enumerate}

若 G 的阶写成 \(p^{r}m\) 的形式，其中 m 不是 p 的倍数，则条件 (a) 与 (b) 等价于 \(\operatorname{Card}(P) = p^{r}\) 。

例子。(1) 在群 \(\mathfrak{S}_p\) 中，设 \(\zeta\) 是一个 \(p\) 阶轮换。由 \(\zeta\) 生成的子群是一个西罗 \(p\) -子群（ \(\mathfrak{S}_p\) 的），因为 \(p\) 不整除 \((p - 1)!\) 。

\begin{enumerate}
\def\labelenumi{(\arabic{enumi})}
\setcounter{enumi}{1}
\tightlist
\item
  *设 \(k\) 是特征为 \(p\) 的有限域，设 \(n\) 为正整数。严格三角群 \(T_1(n, k)\) 是一个西罗 \(p\) -子群（群 \(\mathbf{GL}(n, k)\) 的）。*
\end{enumerate}

定理 2. 每个有限群都包含一个西罗 p-子群。

证明依赖于以下引理。

引理. 设 \(n = p^r m\) ，其中 \(m\) 是一个不是 \(p\) 的倍数的整数。则

\[
\binom {n} {p ^ {r}} \not \equiv 0 \pmod {p}.
\]

设 S 是一个 \(p^r\) 阶群（例如 \(\mathbf{Z} / p^r\mathbf{Z}\) ），T 是一个有 \(m\) 个元素的集合。记 \(\mathbf{X} = \mathbf{S} \times \mathbf{T}\) ，并设 E 是 X 的含 \(p^r\) 个元素的子集所成的集合。则 \(\operatorname{Card}(\mathbf{X}) = n\) ，从而 \(\operatorname{Card}(\mathbf{E}) = \binom{n}{p^r}\) （集合论，III, § 5, 第 8 号，命题 11 的推论 1）。设 S 通过 \(s.(x,y) = (sx,y)(s,x \in \mathbf{S},y \in \mathbf{T})\) 作用在 X 上，并考虑该作用到 E 上的典范扩张。在第 5 号命题 11 的记号中，集合 \(\mathbf{E}^{\mathbf{S}}\) 由 X 的轨道组成，即由子集 \(Y \subset X\) （形如 \(S \times \{t\}, t \in T\) ）所成的集合，因此 \(\operatorname{Card}(\mathbf{E}^{\mathbf{S}}) = m\) 。由第 5 号命题 11，

\[
\binom {n} {p ^ {r}} = \operatorname{Card} (\mathrm{E}) \equiv \operatorname{Card} (\mathrm{E} ^ {\mathrm{S}}) = m \not \equiv 0 \pmod {p},
\]

这便证明了引理。

现在我们证明该定理。设 G 为有限群，n 为其阶；我们记 \(n = p^{r}m\) ，其中 m 不是 p 的倍数。设 E 为 G 的含 \(p^{r}\) 个元素的子集所成的集合。则

\[
\operatorname{Card} (\mathbf {E}) = \binom{n}{p ^ {r}};
\]

从而由引理， \(\operatorname{Card}(\mathbf{E}) \not\equiv 0 (\operatorname{mod.} p)\) 。考虑 G 在自身上按左平移的作用到 E 上的扩张。存在 \(X \in E\) ，其轨道的基数模 p 非零。若 \(H_{X}\) 表示 X 的稳定化子，则 \((G:H_{X}) \not\equiv 0 (\text{mod. } p)\) ，这意味着 \(p^{r}\) 整除 \(\text{Card}(H_{X})\) 。但 \(H_{X}\) 由 \(s \in G\) （满足 \(sX = X\) ）组成；若 \(x \in X\) ，则 \(H_{X} \subset X.x^{-1}\) ，从而 \(\text{Card}(H_{X}) \leqslant \text{Card}(X) = p^{r}\) 。因此 \(\text{Card}(H_{X}) = p^{r}\) 。

推论. 若 p 整除 G 的阶，则群 G 包含一个 p 阶元素。

由定理 2，这归结为 G 是一个 p-群 \(\neq\{e\}\) 的情形；若 \(x\in G\) 不同于 e，则由 x 生成的循环群的阶为 \(p^{n}\) （ \(n\geqslant1\) ），因此包含一个 p 阶子群。

注记. 对每个整除 \(\operatorname{Card}(\mathbf{G})\) 的素数 q，设 \(P_{q}\) 为 G 的一个西罗 q-子群。则 G 中由诸 \(P_{q}\) 生成的子群 H 的阶对每个 q 都是 \(\operatorname{Card}(\mathbf{P}_{q})\) 的倍数，同时又是 \(\operatorname{Card}(\mathbf{G})\) 的约数，因此它等于 G。

定理 3. 设 G 是一个有限群。

\begin{enumerate}
\def\labelenumi{(\alph{enumi})}
\item
  西罗 \(p\) -子群（ \(\mathbf{G}\) 的）彼此共轭。它们的数目模 \(p\) 与 1 同余。
\item
  G 的每个 p-子群都包含在一个西罗 p-子群中。
\end{enumerate}

设 P 是 G 的一个西罗 p-子群（定理 2），并设 H 是 G 的一个 p-子群。令 E = G/P，并考虑 H 在 G/P 上的作用。由于 \(\text{Card}(E) \not\equiv 0 \mod p\) ，第 5 号的命题 11 表明，存在 \(x \in G/P\) 使得对所有 \(h \in H\) 有 hx = x。若 g 是 x 在 G 中的一个代表元，这就意味着 \(H \subset gPg^{-1}\) ，由此得出断言 (b)。

若 H 是一个西罗 p-子群，则 \(\operatorname{Card}(\mathrm{H}) = \operatorname{Card}(\mathrm{P}) = \operatorname{Card}(g\mathrm{Pg}^{-1})\) ，因此 \(H = gPg^{-1}\) ，这证明了 (a) 的第一个断言。

我们现在证明 (a) 的第二个断言。设 \(\mathcal{S}\) 为 G 的西罗 \(p\) -子群所成的集合，并设 P 通过内自同构作用在 \(\mathcal{S}\) 上。元素 \(\mathrm{P} \in \mathcal{S}\) 在此作用下是不动点，我们证明它是唯一的不动点。设 \(\mathrm{Q} \in \mathcal{S}\) 为一个不动点；Q 是被 P 正规化的 G 的西罗子群，因此 P 包含于 Q 的正规化子 N 中。群 P 和 Q 是 N 的西罗 \(p\) -子群；因此存在 \(n \in N\) 使得 \(\mathrm{P} = n\mathrm{Q}n^{-1} = \mathrm{Q}\) 。由第 5 号命题 11， \(\operatorname{Card}(\mathcal{S}) \equiv \operatorname{Card}(\mathcal{S}^{\mathrm{P}}) = 1 (\text{mod. } p)\) 。

推论 1. 设 P 是 G 的一个西罗 p-子群，设 N 是其在 G 中的正规化子，并设 M 是 G 的一个包含 N 的子群。则 M 在 G 中的正规化子等于 M。

设 \(s \in G\) 满足 \(sMs^{-1} = M\) 。子群 \(sPs^{-1}\) （M 的）是 M 的一个西罗 \(p\) -子群。因此存在 \(t \in M\) 使得 \(sPs^{-1} = tPt^{-1}\) ；于是 \(t^{-1}s \in N\) ，从而 \(s \in tN \subset M\) 。

推论 2. 设 \(f: \mathbf{G}_{1} \to \mathbf{G}_{2}\) 是有限群之间的同态。对每个西罗 \(p\) -子群 \(\mathbf{P}_{1}\) （ \(\mathbf{G}_{1}\) 的），存在一个西罗 \(p\) -子群 \(\mathbf{P}_{2}\) （ \(\mathbf{G}_{2}\) 的），使得 \(f(\mathbf{P}_{1}) \subset \mathbf{P}_{2}\) 。

这由定理 3 (b) 应用于子群 \(f(\mathbf{P}_1)\) （ \(\mathbf{G}_2\) 的）得出。

推论 3. (a) 设 H 是 G 的子群。对 H 的每个西罗 p-子群 P，存在 G 的一个西罗 p-子群 Q，使得 P = Q ∩ H。

\begin{enumerate}
\def\labelenumi{(\alph{enumi})}
\setcounter{enumi}{1}
\item
  反之，若 \(\mathbf{Q}\) 是一个西罗 \(p\) -子群（ \(\mathbf{G}\) 的），且 \(\mathbf{H}\) 在 \(\mathbf{G}\) 中正规，则群 \(\mathbf{Q} \cap \mathbf{H}\) 是一个西罗 \(p\) -子群（ \(\mathbf{H}\) 的）。
\item
  (a) \(p\) -群 \(\mathbf{P}\) 包含于一个西罗 \(p\) -子群 \(\mathbf{Q}\) （ \(\mathbf{G}\) 的），且 \(\mathbf{Q} \cap \mathbf{H}\) 是一个 \(p\) -子群（ \(\mathbf{H}\) 的），包含 \(\mathbf{P}\) ，从而等于 \(\mathbf{P}\) 。
\item
  (b) 设 \(\mathbf{P}'\) 是 H 的一个西罗 \(p\) -子群。存在元素 \(g \in G\) 使得 \(g\mathbf{P}'g^{-1} \subset \mathbf{Q}\) 。由于 H 正规， \(\mathbf{P} = g\mathbf{P}'g^{-1}\) 包含于 H 中，从而包含于 \(\mathbf{Q} \cap \mathbf{H}\) 。由于 \(\mathbf{Q} \cap \mathbf{H}\) 是 H 的一个 \(p\) -子群，而 P 是 H 的一个西罗 \(p\) -子群，故 \(\mathbf{P} = \mathbf{Q} \cap \mathbf{H}\) 。
\end{enumerate}

推论 4. 设 N 是 G 的一个正规子群。G 的一个西罗 p-子群在 G/N 中的像是 G/N 的一个西罗 p-子群，且 G/N 的每个西罗 p-子群都可以由此得到。

设 \(G' = G/N\) ，并设 \(P'\) 为 \(G'\) 中的像，来自西罗 \(p\) -子群 \(P\) （ \(G\) 的）。群 \(G\) 传递地作用在 \(G'/P'\) 上，因此 \(G'/P'\) 与 \(G/S\) 等势，其中 \(S\) 是 \(G\) 的一个包含 \(P\) 的子群。于是 \((G':P')\) 整除 \((G:P)\) ，从而不是 \(p\) 的倍数，故 \(p\) -群 \(P'\) 是一个西罗 \(p\) -子群（ \(G'\) 的）。设 \(Q'\) 是另一个西罗 \(p\) -子群（ \(G'\) 的）；则 \(Q' = g'P'g'^{-1}\) 对某个 \(g' \in G'\) 成立；若 \(g \in G\) 是 \(g'\) 的一个代表元，则群 \(Q'\) 是 \(Q = gPg^{-1}\) 的像。

